\section{Classification and computational proof}\label{sec:computation}

\subsection{Physical equivalence and graph certificates}\label{sec:classification}

The equivalence group on \(\ell\) Pauli sites is
\(G_\ell=S_\ell\ltimes\mathrm{GL}(2,2)^\ell\), of order
\(\ell!6^\ell\). A local invertible map permutes the three nonzero
letters, and over \(\F\) every such permutation is symplectic.
It preserves physical weight, commutation, hull dimension, and coset
distance. One representative per orbit is enough.
Two codes in the same orbit differ only by a relabeling of coordinates
and Pauli letters, so an extension exists for one if and only if a
correspondingly relabeled extension exists for the other. An orbit is
one equivalence class, and \(\Aut C\) is the subgroup fixing a code
\(C\). Orbit--stabilizer gives \(|G_\ell|/|\Aut C|\)
distinct labeled copies of that code.
Labeled means that the coordinate positions and the three nonidentity
Pauli letters at each coordinate are fixed.

One graph representation has a vertex for each codeword,
three letter vertices per coordinate, and a coordinate vertex joined to
its three letters. Each codeword vertex is joined to the letter it uses
at each nonzero coordinate. Codeword, letter, and coordinate vertices
have separate colors.
A color-preserving graph isomorphism gives exactly a site permutation and
local letter permutations carrying the code to the other code.
Conversely, every physical equivalence induces a graph isomorphism.
The permutation on word vertices is forced by their incidences, so there
is no additional graph-automorphism kernel. Canonical graph certificates
and automorphism orders are computed with nauty \cite{nauty}.
Here a canonical graph certificate is a representation that is the same
exactly for isomorphic colored graphs. It recognizes duplicate code
classes; it is different from an infeasibility certificate for an
enumerator system.
The original census uses a second encoding with letter triangles rather
than coordinate vertices. The two encodings are compared using the same code generators.

\subsection{Completeness by extension and incidence}

We use both shortening and puncturing. Shortening restricts to the
words that are zero in a coordinate before deleting it; puncturing
deletes the coordinate from every word.

For lengthening a distance-\(d\) code by one coordinate, let the new
coordinate image have binary rank \(\rho\in\{0,1,2\}\). Its zero-symbol
shortening is a parent \(P\) whose dimension is lower by \(\rho\).
The extension is specified by a \(\rho\)-space in the old Pauli quotient
modulo \(P\), with every nonzero coset minimum at least \(d-1\).
The new nonzero symbol adds one unit of weight. Rank one has three
possible nonzero output letters; rank two has six output isomorphisms.
These exhaust all children, because every child has a unique shortening
at the last physical coordinate.
The reason for counting labeled children as well as equivalence classes
is completeness. Producing pairwise inequivalent codes gives a lower
bound on the number of classes. A complete parent list lets us count
every possible child before identifying equivalences. If the listed
child orbits account for that entire count, none is missing.

Let \(\mathcal L_\rho(P)\) be the set of \(\rho\)-dimensional subspaces
\(U\leq V_\ell/P\) whose nonzero cosets have minimum weight at least
\(d-1\). Put
\[
 e_0(P)=1,\qquad e_1(P)=3|\mathcal L_1(P)|,\qquad
 e_2(P)=6|\mathcal L_2(P)|.
\]
The factors three and six count the injections \(\beta:U\to\F^2\).
Each pair \((U,\beta)\) defines the child
\(\{(p+\widetilde u,\beta u):p\in P,u\in U\}\), independently
of the chosen representatives \(\widetilde u\).
Conversely, shortening a child at its last coordinate recovers \(P\),
and the remaining cosets recover \((U,\beta)\). Each labeled child
therefore occurs exactly once. For a fixed child dimension \(c\), the
parents in the following sum have dimension \(c-\rho\), and the
children have dimension \(c\). Orbit--stabilizer counting gives
\begin{equation}\label{eq:lengthening-mass}
 \sum_{\rho,P}\frac{\ell!6^\ell}{|\Aut P|}\,e_\rho(P)
 =\sum_D\frac{(\ell+1)!6^{\ell+1}}{|\Aut D|}
\end{equation}
The sum over \(P\) uses one representative of every parent class of
dimension \(c-\rho\) and minimum distance at least \(d\); the sum
over \(D\) uses the listed child classes of dimension \(c\) and
minimum distance at least \(d\).
The left-hand side counts all labeled extensions of the complete parent
lists. The right-hand side counts the disjoint orbits of the listed children. Each child is checked individually for rank, distance, and
distinct canonical certificate. Equality, together with complete parents,
leaves no room for a missing orbit. Completeness follows inductively from the parent lists.

For the sixteen-qubit application, the lengthening induction starts with
the zero code on four coordinates and retains every parent dimension
that can reach the target. The base is complete because a nonzero word
on four coordinates cannot have weight five. The four required classifications have the
following masses \(\sum_C |G_\ell|/|\Aut C|\):
\begin{center}\small
\begin{tabular}{@{}rrrr@{}}\toprule
\(\ell\)&Binary dimension&Classes&Labeled mass\\\midrule
7&4&43&21,776,221,440\\
8&6&1,579&98,005,190,492,160\\
9&8&2,298&7,958,195,288,064,000\\
10&10&37&2,333,060,795,842,560\\\bottomrule
\end{tabular}
\end{center}
The same induction gives no codes in the three larger-dimension sets
used in Lemma~\ref{lem:16-rankone}. The source release includes
\pathcode{verification/release16/census/check_cells.py}, which reconstructs these classifications,
checks the mass identity at every step, and compares the resulting
canonical certificates with the representatives used in the searches.

The seventeen-qubit classifications start from an elementary base.
For a binary dimension-three code, each coordinate map
\(\F^3\to\F^2\) is determined up to local output equivalence by its
row space. There are fifteen row spaces of dimension at most two:
one zero space, seven lines, and seven planes.
Every multiset of seven such types is tested, giving
\(\binom{21}{7}=116280\) possibilities. Exactly 2,171 meet the
dimension and distance-five conditions; their physical graph certificates
give 32 classes. This establishes the complete seven-site dimension-three
classification directly.

Higher-dimensional supercodes supply another incidence check. Each
dimension-\((c+1)\) binary code has \(2^{c+1}-1\) hyperplanes.
Its dimension-\(c\) parent extensions correspond to nonzero quotient
cosets of minimum weight at least \(d\). Counting parent/supercode
incidences gives the following identity, where \(e(P)\) is the number
of those nonzero cosets:
\begin{equation}\label{eq:incidence}
 \sum_P\frac{|G_\ell|}{|\Aut P|}\,e(P)
 =(2^{c+1}-1)\sum_D\frac{|G_\ell|}{|\Aut D|}.
\end{equation}
The repository also checks the equality separately for each parent orbit
by identifying every hyperplane of every listed child. This prevents
incorrect per-class counts from canceling in a grand total.

Using \eqref{eq:incidence}, the 32 dimension-three parents prove the
43-class seven-site dimension-four list complete; those prove the
one-class dimension-five list complete. The three parent dimensions
then prove the 26,299-class eight-site dimension-five distance-five list
complete through \eqref{eq:lengthening-mass}. Its labeled mass is
\(1,593,099,450,132,480\).

Finally, puncturing any nine-site dimension-five distance-six code
preserves dimension and gives one of those eight-site parents: a
nonzero puncture kernel would be a weight-one word. The new coordinate
is a pair of binary linear forms on \(\F^5\). All 1,024 ordered pairs
are tested per parent. A lift is valid if and only if every parent word
of weight five receives a nonzero new symbol. The unrestricted valid
pairs number 1,791,031, and local output equivalence gives 299,175
canonical lifts. Their complete physical certificate set contains
32,865 classes, with hull counts \(26752,5995,118\) at hull dimensions
one, three, five.

The nine-site dimension-six classification is independently checked by
\eqref{eq:incidence} from those dimension-five classes. Each child has
63 hyperplanes. The incidence count is
\(4,332,624,025,927,680\); the 56 listed child orbits have total labeled
mass \(68,771,809,935,360\), exactly the former number divided by 63.
All 3,528 hyperplanes are identified, and the equality is checked
pointwise across all 32,865 parents. These checks establish completeness of the classifications used in the
seventeen-qubit extension proof.

\subsection{The local code classifications}\label{sec:local-classifications}

The local centralizer codes used in Section~\ref{sec:seventeen} are
obtained by puncturing and lifting. Puncturing an eight-coordinate code of minimum
distance six preserves its binary dimension and gives a seven-coordinate
code of minimum distance at least five. For dimension four, the 43 parent
classes above suffice. Testing all \(2^8=256\) ordered pairs
of linear forms per parent gives 861 valid lifts and 28 equivalence
classes. Fifteen classes contain a word of weight eight; six of these
are isotropic. Exactly three contain a full-support plane. Writing
\(B\) for each of these three codes, we test all 256 cosets in
\(B^\perp/B\). In total, two cosets have
minimum weight six, and the resulting dimension-five codes form one
equivalence class. This supplies the local code needed in
Proposition~\ref{prop:17-plane}.

For the rank-one case, a code \(H\) of dimension \(q+1\) is marked
by a full-support word \(p\in\rad(H)\). At \(q=1\), normalize
\(p=Z^8\) and test every possible second generator, modulo addition
of \(p\). At \(q=2,3\), lift every seven-coordinate parent of dimension
three or four and mark every full-support radical word. The complete
parent lists have 32 and 43 classes. The marked physical graph
certificates give, respectively, four types at \((q,t)=(1,0)\),
six and two at \((2,0),(2,2)\), and seven and five at
\((3,0),(3,2)\). Marking \(p\) means giving its word vertex a
separate color. An isomorphism must preserve the selected word.
The repository includes a fresh coordinate-lift checker for
the full-plane classifications; the complete marked-code check is
identified in Appendix~\ref{app:evidence}.

\subsection{Extension computations}

For each shortened code, the search computes all coset minimum weights
and enumerates the extension maps in the domain established by the
preceding reductions. In the seventeen-qubit cases, it first assigns the
images of a totally isotropic half of the local quotient and then solves
for every complementary image. The Gram equation determines part of the
latter assignment; all remaining entries, including diagonal entries,
are searched. Appendix~\ref{app:search} gives the equations, pruning
conditions, and symmetry arguments.

The computation is divided by shortened-code class, local kernel,
radical identification, and first syndrome image. A case is excluded
only when every part has been exhausted and no extension satisfies the
required equations and weights. Interrupted computations are rerun to
completion. The verification programs reconstruct this decomposition
and check the completion records for every expected part.

\subsection{Implementation checks}

Known \(\qc{5,1,3}\), \(\qc{16,4,5}\), and \(\qc{16,1,6}\)
codes provide positive tests. Seventeen-coordinate \(\qc{17,3,3}\)
examples additionally test prescribed local kernels, shortened-code
kernels, and singular Gram matrices. Their stabilizers and centralizers
are constructed explicitly; ranks, commutation, and minimum distance are
checked by enumeration.

The complementary-image algorithm is compared with direct enumeration
on small examples, including assignments with nonzero diagonal entries.
For the mixed rank-two cases, these comparisons cover 189,860 tuples
in 1,804 fixed-tuple and distance-threshold instances. Separate C
implementations are compared on common inputs. Tests with deliberately
incomplete or inconsistent records check that these cannot be accepted
as complete exclusions. The reproducibility appendix gives the precise
scope of each verification command.
